\documentclass{amsart}

% 数学相关宏包
\usepackage{amsmath,mathtools,amsthm,amssymb}
\usepackage{bbm}
\usepackage{mathrsfs}
\usepackage{centernot}

% 图形和绘图
\usepackage{graphicx,float}
\usepackage{tikz,tikz-cd}
\usepackage{pgfplots}
\usepackage[framemethod=TikZ]{mdframed}
\usetikzlibrary{decorations.pathmorphing,positioning,arrows,automata,calc,shapes,patterns}
\pgfplotsset{compat=1.18}
\usepgfplotslibrary{statistics}

% 表格和列表
\usepackage{booktabs,multirow,colortbl}
\usepackage{enumitem}
\setlist[itemize,enumerate,description]{noitemsep}
\setlistdepth{9}

% 算法
\usepackage{algorithm,algorithmic}

% 字体和颜色
\usepackage{xcolor,listings}
\usepackage{xeCJK}
\setCJKmainfont{Noto Serif CJK SC}

% 页面和链接
\usepackage{fancyhdr,caption}
\usepackage{hyperref}
\usepackage{cleveref}

% 工具宏包
\usepackage{etoolbox,letltxmacro,docmute}
\usepackage{gensymb}

% 自定义设置
\renewcommand{\arraystretch}{1.5}
\let\originalmathbb\mathbb
\renewcommand{\mathbb}[1]{%
    \ifstrequal{#1}{1}{\mathbbm{#1}}{\originalmathbb{#1}}%
}

% 颜色定义
\definecolor{lightblue}{RGB}{173,216,230}
\definecolor{lightgreen}{RGB}{144,238,144}
\definecolor{lightyellow}{RGB}{255,255,224}
\definecolor{lightgray}{RGB}{211,211,211}

		% 超链接设置
		\hypersetup{
		    colorlinks=true,
		    linkcolor=red,
		    filecolor=magenta,
		    urlcolor=cyan,
		    pdftitle={\@title},
		    pdfauthor={\@author},
		    pdfsubject={数学笔记},
		    pdfkeywords={数学},
		    bookmarksnumbered=true,
		    bookmarksopen=true,
		    bookmarksopenlevel=6,
		    pdfstartview=FitH,
		    pdfpagemode=UseOutlines,
		    pdfborder={0 0 0}
		}

% 定理环境
\theoremstyle{plain}
\newtheorem{theorem}{Theorem}[section]
\newtheorem{lemma}[theorem]{Lemma}
\newtheorem{corollary}[theorem]{Corollary}
\newtheorem{proposition}[theorem]{Proposition}

\theoremstyle{definition}
\newtheorem{definition}[theorem]{Definition}
\newtheorem{example}[theorem]{Example}
\newtheorem{exercise}[theorem]{Exercise}
\newtheorem{problem}[theorem]{Problem}

\theoremstyle{remark}
\newtheorem{remark}[theorem]{Remark}
\newtheorem{note}[theorem]{Note}
\newtheorem{observation}[theorem]{Observation}

% 解答环境
\newtheorem*{solution}{Solution}
\newtheorem*{proof*}{Proof}

% 图片路径
\graphicspath{{F:/iCloudDrive/iCloud~md~obsidian/2025-summer/figures/}}

\author{乐绎华, Yue Yihua}
\address{中山大学, Sun Yat-sen University}
\email{yueyh@mail2.sysu.edu.cn}
\date{\today}
\title{Mathematical Notes}

\begin{document}

\maketitle
\tableofcontents

\chapter{Finding Lyapunov Functions for Markov Chains}

\begin{note}
\textbf{Lecture Notes}\textbf{Speaker}: Weijun Xu, Peking University\textbf{Title}: Finding Lyapunov functions\textbf{Abstract}: Existence of invariant measures for Markov processes are often established by finding suitable Lyapunov functions. In practice, constructing such a function is often more of an art than science. We introduce some techniques in finding Lyapunov functions, and demonstrate their usage in one particularly interesting example.
\end{note}
\section{Introduction}

Finding Lyapunov functions is a fundamental technique in the study of Markov processes, yet it remains one of the most challenging aspects of the theory. While the theoretical framework is well-established, the actual construction of these functions often requires creativity, intuition, and experience.

\subsection{Why Lyapunov Functions Matter}

Lyapunov functions serve multiple crucial purposes in Markov chain theory:

\begin{enumerate}
	\item \textbf{Proving Existence of Invariant Measures}: Through the Krylov-Bogoliubov theorem, Lyapunov functions establish tightness of empirical measures, which implies existence of invariant measures.
	\item \textbf{Quantifying Convergence Rates}: Different forms of drift conditions yield different convergence speeds:
	\begin{itemize}
		\item $PV - V \leq -cV$ gives geometric (exponential) convergence
		\item $PV - V \leq -cV^{\beta}$ with $\beta < 1$ gives polynomial convergence
		\item The function $\Phi$ in $PV - V \leq -\Phi(V)$ directly determines the convergence rate
	\end{itemize}
	\item \textbf{Establishing Moment Bounds}: If $\mathbb{E}_\nu[V] < \infty$ for the invariant measure $\nu$, then $V$ provides moment bounds for the stationary distribution.
	\item \textbf{Designing Efficient Algorithms}: In MCMC, Lyapunov functions guide the choice of proposal distributions and adaptation strategies.
\end{enumerate}

\subsection{The Core Challenge}

The fundamental difficulty lies in the conflicting requirements:

\begin{itemize}
	\item \textbf{Too Weak}: If $V$ grows too slowly (e.g., $V(x) = \log(1 + |x|)$), it may not provide useful moment bounds or fast convergence rates.
	\item \textbf{Too Strong}: If $V$ grows too quickly (e.g., $V(x) = e^{|x|^2}$), the drift condition $PV \leq \lambda V + b$ may be impossible to verify.
	\item \textbf{Just Right}: We need $V$ to grow at exactly the right rate to both control the chain and satisfy the drift condition.
\end{itemize}

\subsection{Systematic Approaches}

Rather than relying purely on intuition, we'll explore systematic techniques:

\begin{enumerate}
	\item \textbf{Start with the Chain's Structure}:
	\begin{itemize}
		\item Linear chains → quadratic Lyapunov functions
		\item Polynomial noise → polynomial Lyapunov functions
		\item Exponential tails → exponential Lyapunov functions
	\end{itemize}
	\item \textbf{Work Backwards from Desired Properties}:
	\begin{itemize}
		\item Want geometric ergodicity? Need $PV - V \leq -cV$
		\item Want finite $p$-th moments? Try $V(x) = 1 + |x|^p$
		\item Want to prove recurrence? Any unbounded $V$ with negative drift suffices
	\end{itemize}
	\item \textbf{Use Physical Intuition}:
	\begin{itemize}
		\item Energy functions in physics often make good Lyapunov functions
		\item Distance to equilibrium in optimization problems
		\item Potential functions from gradient flows
	\end{itemize}
	\item \textbf{Exploit Problem-Specific Features}:
	\begin{itemize}
		\item Symmetries can simplify the search
		\item Conservation laws suggest natural candidates
		\item Known invariant measures hint at appropriate growth rates
	\end{itemize}
\end{enumerate}

\subsection{A Concrete Example: The Art in Practice}

Consider a simple nonlinear autoregression: $X_{n+1} = \theta X_n + \sigma(X_n) Z_n$ where $Z_n$ are i.i.d. standard normal.

\textbf{First Attempt}: Try $V(x) = 1 + x^2$

\begin{itemize}
	\item Compute: $PV(x) = 1 + \theta^2 x^2 + \sigma(x)^2$
	\item Need: $\theta^2 x^2 + \sigma(x)^2 \leq \lambda x^2 + b$ for some $\lambda < 1$
\end{itemize}

\textbf{The Art}:

\begin{itemize}
	\item If $\sigma(x) = \sigma_0$ (constant), this works when $|\theta| < 1$
	\item If $\sigma(x) = \sigma_0 \sqrt{1 + x^2}$, we need to modify $V$
	\item If $\sigma(x) = \sigma_0 |x|$, we might need $V(x) = 1 + |x|^p$ for some $p < 2$
\end{itemize}

This interplay between the chain's dynamics and the choice of $V$ exemplifies why finding Lyapunov functions is "more art than science."

\subsection{Roadmap of This Lecture}

We will proceed systematically through:

\begin{enumerate}
	\item \textbf{Theoretical Foundations}: Adjoint operators, Krylov-Bogoliubov theorem, and the connection to supermartingales
	\item \textbf{Basic Techniques}: Quadratic forms, distance-based functions, and the test function method
	\item \textbf{Advanced Methods}: Small sets, polynomial rates, and adaptive MCMC
	\item \textbf{Case Studies}: Detailed examples showing the thought process behind choosing $V$
	\item \textbf{Common Pitfalls}: What goes wrong and how to fix it
\end{enumerate}

By the end, you'll have both theoretical understanding and practical tools for finding Lyapunov functions in your own research.

\section{Setup}

\begin{itemize}
	\item State space $\mathcal{X}$
	\item Feller Markov chain $(X_n)_{n\geq0}$ on $\mathcal{X}$
	\item Markov operator $P:\mathcal{G}(\mathcal{X})\to \mathcal{G}(\mathcal{X})$.
	\item $(Pf)(x)\coloneqq \mathbb{E}(f(X_1)\mid X_0=x)=\mathbb{E}^{x}(f(X_1))$.
\end{itemize}

Question:

\begin{itemize}
	\item Existence of IM (invariant measure)?
	\item Speed of convergence?
\end{itemize}

\section{Adjoint of Markov Operators}

For a Markov operator $P: \mathcal{G}(\mathcal{X}) \to \mathcal{G}(\mathcal{X})$ acting on functions, its adjoint $P^*$ acts on measures and is defined by the duality relation:
\[
\langle Pf, \mu \rangle = \langle f, P^*\mu \rangle
\]

or equivalently:
\[
\int (Pf)(x) \, \mu(dx) = \int f(x) \, (P^*\mu)(dx)
\]

\subsection{Key Properties:}

\begin{itemize}
	\item If $\mu$ is the distribution of $X_0$, then $P^*\mu$ is the distribution of $X_1$
	\item For Dirac measure: $(P^*\delta_x)(A) = P(x,A) = \mathbb{P}^x(X_1 \in A)$
	\item The $k$-th power: $(P^*)^k\mu$ is the distribution of $X_k$ given initial distribution $\mu$
\end{itemize}

\begin{theorem}[krytov-Bogotioubov]
Let $P^{*}$ be the adjoint of $P$ and $(P^{*})^{k}=P^{*}_k$, if $\exists x\in \mathcal{X}$ s.t.  $\mu _n\coloneqq \frac{1}{n}\sum_{k=0}^{n-1}P^{*}_k\delta_{x}$ is tight in $n$ then (Feller) $\Rightarrow$ existence of IM.
\end{theorem}
\begin{definition}[Tight]
A family of probability measures $\left\{\mu_\alpha\right\}_{\alpha \in I}$ on a metric space $S$ (often $\mathbb{R}^d$) is called \textbf{tight} if for every $\varepsilon>0$, there exists a compact set $K_\varepsilon \subset S$ such that for all $\alpha \in I$,
\[
\mu_\alpha(K_\varepsilon) > 1 - \varepsilon.
\]
\end{definition}
\begin{proof}
Tightness implies that $\exists T_m\to +\infty$ s.t. $\nu _m\coloneqq \frac{1}{T_m}\sum_{k=0}^{T_m-1}P^{*}_k\delta_{x}$ converges weakly to $\nu$. Then
\[
\begin{aligned}
P^{*}\nu & =P^{*}(\lim_{ m \to \infty }\nu _m)\overset{ \text{Feller} }{ = }\lim_{ m \to \infty }P^{*}\nu _m=\lim_{ m \to \infty } \frac{1}{T_m}\sum_{k=0}^{T_m-1} P^{*}_{k+1}\delta_{x} \\
 & =\lim_{ m \to \infty } \frac{1}{T_m}\sum_{k=1}^{T_m} P^{*}_{k}\delta_{x} \\
 & = \lim_{ m \to \infty } \frac{1}{T_m}\left(\sum_{k=0}^{T_m-1} P^{*}_{k}\delta_{x} + P^{*}_{T_m}\delta_{x} - P^{*}_0\delta_{x}\right) \\
 & = \lim_{ m \to \infty } \left(\nu_m + \frac{1}{T_m}P^{*}_{T_m}\delta_{x} - \frac{1}{T_m}\delta_{x}\right)
\end{aligned}
\]
Since $P^{*}_{T_m}\delta_{x}$ is a probability measure and $T_m \to +\infty$, we have:
\[
\frac{1}{T_m}P^{*}_{T_m}\delta_{x} \to 0 \text{ and } \frac{1}{T_m}\delta_{x} \to 0
\]

Therefore:
\[
P^{*}\nu = \lim_{ m \to \infty } \nu_m = \nu
\]

This shows that $\nu$ is an invariant measure (IM) for the Markov chain.
\end{proof}

\section{Lyapunov Functions and Convergence}

Once we have established the existence of an invariant measure, the next question is about the speed of convergence. Lyapunov functions provide a powerful tool for this analysis.

\begin{definition}[Lyapunov Function]
A function $V: \mathcal{X} \to [0,+\infty]$ is called a \textbf{Lyapunov function} for the Markov operator $P$ if:
	\begin{enumerate}
		\item $V$ is lower semi-continuous
		\item The sublevel sets $\{x: V(x) \leq r\}$ are compact for all $r \geq 0$
		\item There exist constants $\lambda \in (0,1)$ and $b < \infty$ such that:
\[
PV(x) \leq \lambda V(x) + b
\]
	\end{enumerate}
\end{definition}
\subsection{Foster-Lyapunov Criteria}

The existence of a Lyapunov function implies:

\begin{itemize}
	\item \textbf{Geometric ergodicity}: If $V$ is bounded on compact sets, then the chain converges exponentially fast to its invariant measure
	\item \textbf{Moment bounds}: $\mathbb{E}_\nu[V] < \infty$ where $\nu$ is the invariant measure
\end{itemize}

\subsection{Connection to Tightness}

The Lyapunov function provides a direct way to verify tightness:

\begin{itemize}
	\item If $PV \leq \lambda V + b$ with $\lambda < 1$, then the empirical measures $\mu_n = \frac{1}{n}\sum_{k=0}^{n-1}P^*_k\delta_x$ are tight
	\item This is because the sublevel sets of $V$ provide the required compact sets for the tightness condition
\end{itemize}

\subsection{Lyapunov Functions and Supermartingales}

Let $V: \mathcal{X} \to [1,+\infty)$ be a Lyapunov function. The drift condition has a deep connection to martingale theory.

\begin{theorem}[Foster's Theorem (Supermartingale Characterization)]
If there exists $\Phi: [1,+\infty) \to \mathbb{R}$ with $\Phi(v) \to +\infty$ as $v \to +\infty$, and the drift condition:
\[
PV(x) - V(x) \leq -\Phi(V(x))
\]
Then the process
\[
M_n = V(X_n) + \sum_{k=0}^{n-1} \Phi(V(X_k))
\]
is a supermartingale with respect to the natural filtration.
\end{theorem}
\textbf{Proof}: We need to show $\mathbb{E}[M_{n+1} \mid \mathcal{F}_n] \leq M_n$. We have:
\[
\mathbb{E}[M_{n+1} \mid \mathcal{F}_n] = \mathbb{E}[V(X_{n+1}) \mid \mathcal{F}_n] + \sum_{k=0}^{n} \Phi(V(X_k))
\]

Since $\mathbb{E}[V(X_{n+1}) \mid \mathcal{F}_n] = PV(X_n)$:
\[
= PV(X_n) + \sum_{k=0}^{n} \Phi(V(X_k))
\]

Using the drift condition $PV(X_n) \leq V(X_n) - \Phi(V(X_n))$:
\[
\leq V(X_n) - \Phi(V(X_n)) + \sum_{k=0}^{n} \Phi(V(X_k)) = V(X_n) + \sum_{k=0}^{n-1} \Phi(V(X_k)) = M_n
\]

Therefore $\mathbb{E}[M_{n+1} \mid \mathcal{F}_n] \leq M_n$, making $(M_n)$ a supermartingale.

\subsection{Applications}

\begin{enumerate}
	\item \textbf{Moment Bounds}: Since $(M_n)$ is a supermartingale and $M_n \geq V(X_n)$, we have $\mathbb{E}^x[V(X_n)] \leq V(x)$ when $\Phi(v) \geq 0$.
	\item \textbf{Recurrence}: If $\Phi(v) > 0$ for all $v > v_0$ for some $v_0$, then the chain returns to the set $\{x: V(x) \leq v_0\}$ in finite expected time.
	\item \textbf{Geometric Ergodicity}: If $\Phi(v) = cv$ for some $c > 0$ (i.e., $PV - V \leq -cV$), then the chain has geometric convergence rate.
	\item \textbf{Tail Bounds}: The supermartingale property combined with concentration inequalities gives exponential tail bounds on hitting times.
\end{enumerate}

\subsection{Example: Classical Drift Condition}

For the standard drift condition $PV \leq \lambda V + b$, we can rewrite:
\[
PV(x) - V(x) \leq (\lambda - 1)V(x) + b
\]

When $V(x) > \frac{b}{1-\lambda}$, we have $PV(x) - V(x) < 0$. More generally, if we define:
\[
\Phi(v) = (1-\lambda)v - b
\]

Then for $v$ large enough, $\Phi(v) > 0$ and the supermartingale becomes:
\[
M_n = V(X_n) + \sum_{k=0}^{n-1}\max\{0, (1-\lambda)V(X_k) - b\}
\]

\section{Methods for Finding Lyapunov Functions}

\subsection{1. Quadratic Forms (for Linear Systems)}

For linear Markov chains on $\mathbb{R}^d$, try $V(x) = x^T Q x$ where $Q$ is positive definite:

\begin{itemize}
	\item Solve the Lyapunov equation: $A^T Q A - Q = -I$ where $A$ is the transition matrix
	\item If solvable with $Q > 0$, then $V(x) = x^T Q x$ is a Lyapunov function
\end{itemize}

\subsection{2. Distance-Based Functions}

Common choices include:

\begin{itemize}
	\item $V(x) = 1 + |x|^p$ for $p \geq 1$
	\item $V(x) = e^{\alpha |x|}$ for some $\alpha > 0$
	\item $V(x) = \log(1 + |x|^2)$
\end{itemize}

\subsection{3. Test Functions Method}

Start with a candidate $V$ and verify:
\[
PV(x) - V(x) = \mathbb{E}^x[V(X_1)] - V(x)
\]

If this is negative for large $|x|$, adjust parameters to achieve the drift condition.

\subsection{4. Coupling Method}

If you can couple the chain with a simpler process:

\begin{itemize}
	\item Find Lyapunov function for the simpler process
	\item Transfer to the original chain via the coupling
\end{itemize}

\section{Convergence Rates from Lyapunov Functions}

\begin{theorem}[Geometric Ergodicity]
If $V \geq 1$ is a Lyapunov function with $PV \leq \lambda V + b$, and $V$ is bounded on compact sets, then:
\[
\|P^n(x,\cdot) - \nu(\cdot)\|_{TV} \leq C V(x) \rho^n
\]
for some $C < \infty$ and $\rho \in (\lambda, 1)$.
\end{theorem}
\subsection{Proof Sketch:}

\begin{enumerate}
	\item The drift condition implies $\mathbb{E}^x[V(X_n)] \leq \lambda^n V(x) + \frac{b}{1-\lambda}$
	\item This gives exponential moment bounds
	\item Apply coupling arguments to get TV convergence
\end{enumerate}

\section{Examples}

\subsection{Example 1: Random Walk with Drift}

Consider $X_{n+1} = \theta X_n + Z_n$ where $|\theta| < 1$ and $Z_n$ are i.i.d. bounded.

\textbf{Lyapunov function}: $V(x) = 1 + x^2$

\textbf{Verification}:
\[
PV(x) = 1 + \mathbb{E}[(θx + Z)^2] = 1 + θ^2 x^2 + \mathbb{E}[Z^2] \leq θ^2 V(x) + (1 + \mathbb{E}[Z^2])
\]

Since $θ^2 < 1$, this gives geometric ergodicity.

\subsection{Example 2: Metropolis-Hastings Algorithm}

For target density $\pi(x) \propto e^{-U(x)}$ with $U(x) \to \infty$ as $|x| \to \infty$:

\textbf{Lyapunov function}: $V(x) = e^{\delta U(x)}$ for small $\delta > 0$

Under appropriate conditions on the proposal, one can show $PV \leq \lambda V + b$.

\subsection{Underdamped Langevin Dynamics}

Consider the second-order Langevin equation (also known as kinetic Langevin or underdamped Langevin):
\[
\begin{cases}
dX_{t}=v_{t}dt \\
dv_{t}=-W'(X_{t})dt-v_{t}dt+\Gamma_2dB_{t} 
\end{cases}
\]
This system describes:

\begin{itemize}
	\item $X_t$: position of a particle
	\item $v_t$: velocity of the particle
	\item $W(x)$: potential energy function
	\item $-v_t dt$: friction/damping term
	\item $\Gamma_2 dB_t$: random thermal noise
\end{itemize}

Then
\[
\mathcal{L}=\partial^{2}_{v}-v\partial_{v}-W'(x)\partial_{v}+v\partial_{x}
\]
\textbf{Derivation}: To find the generator, we apply Itô's formula to a test function $f(x,v)$. For the system:
\[
\begin{cases}
dX_{t}=v_{t}dt \\
dv_{t}=-W'(X_{t})dt-v_{t}dt+\Gamma_2dB_{t} 
\end{cases}
\]
We have:
\[
df(X_t, v_t) = \frac{\partial f}{\partial x} dX_t + \frac{\partial f}{\partial v} dv_t + \frac{1}{2} \frac{\partial^2 f}{\partial v^2} (dv_t)^2
\]

Substituting the dynamics:

\begin{itemize}
	\item $\frac{\partial f}{\partial x} dX_t = v \frac{\partial f}{\partial x} dt$
	\item $\frac{\partial f}{\partial v} dv_t = \frac{\partial f}{\partial v}[-W'(x) - v]dt + \Gamma_2 \frac{\partial f}{\partial v} dB_t$
	\item $(dv_t)^2 = \Gamma_2^2 dt$ (using $(dB_t)^2 = dt$)
\end{itemize}

Therefore:
\[
df = \left[v \frac{\partial f}{\partial x} - W'(x)\frac{\partial f}{\partial v} - v\frac{\partial f}{\partial v} + \frac{\Gamma_2^2}{2}\frac{\partial^2 f}{\partial v^2}\right]dt + \Gamma_2 \frac{\partial f}{\partial v} dB_t
\]

The generator is the drift coefficient (assuming $\Gamma_2 = \sqrt{2}$ for standard form):
\[
\mathcal{L}f = v\partial_x f - W'(x)\partial_v f - v\partial_v f + \partial_{vv}^2 f
\]

This gives us the generator: $\mathcal{L} = v\partial_x - W'(x)\partial_v - v\partial_v + \partial_{vv}^2$.

\subsection{Finding a Lyapunov Function}

For this system, a natural candidate is the total energy (Hamiltonian):
\[
H(x,v) = \frac{1}{2}v^2 + W(x)
\]

Let's check if this satisfies a Lyapunov condition. Computing $\mathcal{L}H$:
\[
\begin{aligned}
\mathcal{L}H &= v\partial_x H - W'(x)\partial_v H - v\partial_v H + \partial_{vv}^2 H \\
&= v \cdot W'(x) - W'(x) \cdot v - v \cdot v + 0 \\
&= -v^2
\end{aligned}
\]

So we have $\mathcal{L}H = -v^2 \leq 0$. This shows energy dissipation but isn't quite a Lyapunov function yet.

\textbf{Modified Lyapunov Function}: Try
\[
V(x,v) = H(x,v) + \epsilon xv = \frac{1}{2}v^2 + W(x) + \epsilon xv
\]

for small $\epsilon > 0$. Computing:
\[
\mathcal{L}V = -v^2 + \epsilon(v^2 - xW'(x) - xv)
\]

Under appropriate growth conditions on $W$ (e.g., $W(x) \geq c_1|x|^2 - c_2$ and $|W'(x)| \leq c_3(1 + |x|)$), we can show:
\[
\mathcal{L}V \leq -c V + d
\]

for some constants $c, d > 0$, giving us a proper Lyapunov function.

\textbf{Key Insight}: The cross term $\epsilon xv$ couples position and velocity to create the necessary drift. This is a common technique for second-order systems.

\section{Practical Considerations}

\subsection{Choosing the Right Lyapunov Function}

\begin{enumerate}
	\item \textbf{Start simple}: Try $V(x) = 1 + |x|^2$ first
	\item \textbf{Match the tail behavior}: If the chain has polynomial tails, use polynomial $V$; if exponential tails, use exponential $V$
	\item \textbf{Use problem structure}: Exploit symmetries or conservation properties
\end{enumerate}

\subsection{Verifying the Drift Condition}

To check $PV(x) \leq \lambda V(x) + b$:

\begin{enumerate}
	\item \textbf{Compute $PV(x)$ explicitly} when possible
	\item \textbf{Use bounds}: Often easier to bound $PV(x)$ than compute exactly
	\item \textbf{Split the state space}: Different bounds for $|x| \leq R$ and $|x| > R$
\end{enumerate}

\subsection{Common Pitfalls}

\begin{itemize}
	\item \textbf{Too weak $V$}: If $V$ grows too slowly, may not get moment bounds
	\item \textbf{Too strong $V$}: If $V$ grows too fast, drift condition may fail
	\item \textbf{Forgetting the compact set condition}: Need sublevel sets to be compact
\end{itemize}

\section{Advanced Topics}

\subsection{Small Sets and Drift Conditions}

\begin{definition}[Small Set]
A set $C \subseteq \mathcal{X}$ is called a \textbf{small set} if there exists $n \geq 1$, $\epsilon > 0$, and a probability measure $\nu$ such that:
\[
P^n(x,A) \geq \epsilon \nu(A) \quad \forall x \in C, A \in \mathcal{B}(\mathcal{X})
\]
\end{definition}
\textbf{Drift and Minorization}: Combine Lyapunov functions with small sets:

\begin{itemize}
	\item Drift condition: $PV \leq \lambda V + b \mathbf{1}_C$
	\item Minorization on $C$: Lower bound on transition probabilities
\end{itemize}

\subsection{Polynomial Convergence}

When geometric ergodicity fails, look for polynomial rates:

\textbf{Polynomial Lyapunov function}: $V(x) = 1 + |x|^{\alpha}$ for $\alpha < 1$

\textbf{Drift condition}: $PV(x) \leq V(x) - c V(x)^{\beta} + b \mathbf{1}_C$ for $\beta < 1$

This gives convergence rate $O(n^{-r})$ for some $r > 0$.

\subsection{Adaptive MCMC}

For adaptive MCMC algorithms, Lyapunov functions help prove:

\begin{itemize}
	\item Containment: The adaptation doesn't destabilize the chain
	\item Diminishing adaptation: The adaptation effect vanishes
\end{itemize}

\subsection{Lyapunov Functions for Continuous-Time Chains}

For continuous-time Markov processes with generator $\mathcal{L}$:

\textbf{Lyapunov condition}: $\mathcal{L}V(x) \leq -c V(x) + d$

This is the continuous analog of the discrete drift condition.

\section{Summary}

Lyapunov functions are fundamental tools for:

\begin{enumerate}
	\item \textbf{Proving existence} of invariant measures (via tightness)
	\item \textbf{Establishing convergence rates} (geometric or polynomial)
	\item \textbf{Obtaining moment bounds} for the invariant measure
	\item \textbf{Designing efficient MCMC algorithms}
\end{enumerate}

The key is finding the right balance: $V$ should grow fast enough to control the chain but slowly enough to satisfy the drift condition.

\section{References}

\begin{enumerate}
	\item \textbf{Meyn, S. P., \& Tweedie, R. L. (2009)}. \textit{Markov Chains and Stochastic Stability} (2nd ed.). Cambridge University Press.
	\begin{itemize}
		\item The definitive reference for Lyapunov function methods in Markov chain theory
	\end{itemize}
	\item \textbf{Hairer, M., \& Mattingly, J. C. (2011)}. Yet another look at Harris' ergodic theorem for Markov chains. \textit{Seminar on Stochastic Analysis, Random Fields and Applications VI}, 109-117.
	\begin{itemize}
		\item Modern treatment of ergodicity using Lyapunov functions
	\end{itemize}
	\item \textbf{Douc, R., Fort, G., Moulines, E., \& Soulier, P. (2004)}. Practical drift conditions for subgeometric rates of convergence. \textit{The Annals of Applied Probability}, 14(3), 1353-1377.
	\begin{itemize}
		\item Extends Lyapunov methods to polynomial convergence rates
	\end{itemize}
	\item \textbf{Roberts, G. O., \& Rosenthal, J. S. (2004)}. General state space Markov chains and MCMC algorithms. \textit{Probability Surveys}, 1, 20-71.
	\begin{itemize}
		\item Excellent survey with practical MCMC applications
	\end{itemize}
	\item \textbf{Fort, G., \& Moulines, E. (2003)}. Polynomial ergodicity of Markov transition kernels. \textit{Stochastic Processes and their Applications}, 103(1), 57-99.
	\begin{itemize}
		\item Comprehensive treatment of polynomial convergence
	\end{itemize}
\end{enumerate}


\end{document}